% Numero di pagina e testatina nel margine esterno, piu' note a margine:
% il margine deve lasciare un alone libero attorno a tutto cio' che e' stampato
\documentclass[a4paper,11pt,twoside]{article}
\usepackage[a4paper,inner=3cm,outer=4.4cm,top=3.6cm,bottom=3.6cm,marginparwidth=2.6cm]{geometry}
\usepackage{lipsum}
\usepackage{fancyhdr}
\pagestyle{fancy}
\fancyhf{}
\fancyhfoffset[LE,RO]{2.4cm}% lato esterno ([O] da solo vorrebbe dire "pagine dispari")
\fancyhead[LE,RO]{\textbf{\thepage}}
\fancyhead[RE,LO]{\small\itshape Layout di un chip come cornice}
\fancyfoot[C]{\footnotesize asicmargins, prova dell'alone}
\usepackage[seed=9,halo=6]{asicmargins}
\begin{document}
\section*{Alone attorno a testatine e note}
\lipsum[1]\marginpar{\footnotesize Nota a margine: il routing deve girarle attorno.}
\lipsum[2-4]\marginpar{\footnotesize Seconda nota, piu' in basso.}
\lipsum[5-9]\marginpar{\footnotesize Nota sulla pagina pari.}
\lipsum[10-20]
\end{document}
